% Created 2025-12-27 sáb 14:04
% Intended LaTeX compiler: pdflatex
\documentclass[11pt]{article}
\usepackage[utf8]{inputenc}
\usepackage[T1]{fontenc}
\usepackage{graphicx}
\usepackage{longtable}
\usepackage{wrapfig}
\usepackage{rotating}
\usepackage[normalem]{ulem}
\usepackage{amsmath}
\usepackage{amssymb}
\usepackage{capt-of}
\usepackage{hyperref}
\usepackage{minted}

\usepackage{parskip}
\usepackage[utf8]{inputenc} %% For unicode chars
\usepackage{placeins}
\usepackage[margin=2.50cm]{geometry}
\usepackage[T1]{fontenc}
\usepackage{mathpazo}
\linespread{1.05}
\usepackage[scaled]{helvet}
\usepackage{courier}
\hypersetup{colorlinks=true,linkcolor=blue}
\RequirePackage{fancyvrb}
\DefineVerbatimEnvironment{verbatim}{Verbatim}{fontsize=\small,formatcom = {\color[rgb]{0.5,0,0}}}
\usepackage{mdframed}
\BeforeBeginEnvironment{minted}{\begin{mdframed}}
\AfterEndEnvironment{minted}{\end{mdframed}}
\setcounter{secnumdepth}{3}
\date{}
\title{}
\hypersetup{
 pdfauthor={},
 pdftitle={},
 pdfkeywords={},
 pdfsubject={},
 pdfcreator={Emacs 27.1 (Org mode 9.6.18)}, 
 pdflang={Spanish}}
\begin{document}

\setcounter{tocdepth}{3}
\tableofcontents

\setlength\parindent{10pt}

\section{To-Do Class}
\label{sec:org40600f6}

Let's create a comprehensive \texttt{Todo} class for our To-Do
application in PHP. This class will include methods for various CRUD
operations such as creating, reading, updating, and deleting To-Do
items, as well as marking them as done or undone. Here's how you can
structure it:

\begin{minted}[]{php}
<?=
  class Todo {
      private $pdo;

      public function __construct($pdo) {
          $this->pdo = $pdo;
      }

      public function create($userId, $title, $description, $categoryId) {
          $sql = "INSERT INTO todos (user_id, title, description, category_id) VALUES (?, ?, ?, ?)";
          $stmt = $this->pdo->prepare($sql);
          $stmt->execute([$userId, $title, $description, $categoryId]);
      }

      public function read($todoId) {
          $sql = "SELECT * FROM todos WHERE id = ?";
          $stmt = $this->pdo->prepare($sql);
          $stmt->execute([$todoId]);
          return $stmt->fetch();
      }

      public function update($todoId, $title, $description, $categoryId) {
          $sql = "UPDATE todos SET title = ?, description = ?, category_id = ? WHERE id = ?";
          $stmt = $this->pdo->prepare($sql);
          $stmt->execute([$title, $description, $categoryId, $todoId]);
      }

      public function delete($todoId) {
          $sql = "DELETE FROM todos WHERE id = ?";
          $stmt = $this->pdo->prepare($sql);
          $stmt->execute([$todoId]);
      }

      public function markDone($todoId) {
          $sql = "UPDATE todos SET done = TRUE WHERE id = ?";
          $stmt = $this->pdo->prepare($sql);
          $stmt->execute([$todoId]);
      }

      public function markUndone($todoId) {
          $sql = "UPDATE todos SET done = FALSE WHERE id = ?";
          $stmt = $this->pdo->prepare($sql);
          $stmt->execute([$todoId]);
      }

      public function getAllTodosByUser($userId) {
          $sql = "SELECT * FROM todos WHERE user_id = ?";
          $stmt = $this->pdo->prepare($sql);
          $stmt->execute([$userId]);
          return $stmt->fetchAll();
      }
  }
\end{minted}

\subsubsection{Explanation:}
\label{sec:org6892ada}
\begin{itemize}
\item \textbf{Constructor}: The constructor takes a PDO database connection object.
This object is used for all database operations within the class.

\item \textbf{Create Todo}: The \texttt{create} method adds a new todo item to the
database for a specific user.

\item \textbf{Read Todo}: The \texttt{read} method fetches the details of a specific todo
item by its ID.

\item \textbf{Update Todo}: The \texttt{update} method modifies the details of an existing
todo item.

\item \textbf{Delete Todo}: The \texttt{delete} method removes a todo item from the
database.

\item \textbf{Mark Done/Undone}: The \texttt{markDone} and \texttt{markUndone} methods update the
status of a todo item to either completed or not completed.

\item \textbf{Get All Todos}: The \texttt{getAllTodosByUser} method retrieves all todo
items for a specific user.
\end{itemize}

\subsubsection{Usage Example:}
\label{sec:org34b20b6}
Here's an example of how you might use this class in your application:

\begin{minted}[]{php}
<?=
$pdo = new PDO("mysql:host=localhost;dbname=your_database", "username", "password");
$todo = new Todo($pdo);

// Create a new Todo
$todo->create(1, "Buy groceries", "Milk, eggs, bread", 2);

// Update a Todo
$todo->update(1, "Buy groceries and vegetables", "Milk, eggs, bread, and carrots", 2);

// Mark a Todo as done
$todo->markDone(1);

// Get all Todos for a user
$userTodos = $todo->getAllTodosByUser(1);
\end{minted}

\subsubsection{Security Considerations:}
\label{sec:org57b85d8}
\begin{itemize}
\item Always use prepared statements as done in the methods to prevent SQL
Injection attacks.
\item Validate and sanitize all inputs before passing them to these methods.
\item Handle exceptions and errors appropriately, especially for database
operations.
\item Ensure that the user performing the operations has the right
permissions, especially for updating and deleting operations.
\end{itemize}

This \texttt{Todo} class provides a solid foundation for managing todos in our
application with standard CRUD operations and specific functionalities
like marking todos as done or undone.
\end{document}
